\section{预期目标及进度安排}

\subsection{预期成果}
本课题围绕“基于深度学习的三维物体材质生成与迁移”展开，  
预期在算法研究与系统实现两个层面取得如下成果：  

（1）在方法层面，设计并实现一种基于深度学习的三维材质生成与迁移模型。  
该模型能够以图片或者带材质的参考模型与对应白模为输入，  
结合几何结构信息，  
自动生成与目标形态一致的高质量材质贴图，  
实现从白模到完整材质模型的自动化转换。  
同时，  
方法应具备一定的泛化能力，  
能够适用于不同类别的三维物体。  

（2）在效果层面，生成的材质应在视觉上具有较好的真实感与一致性，  
能够合理反映模型的几何结构与细节特征。  
在多视角下保持较好的连续性与一致性，  
减少纹理错位、拉伸或失真等问题。  

（3）在实验与评估方面，构建相应的数据集与实验流程，  
对所提出的方法进行定性与定量评估。  
通过与传统方法或已有方法对比，  
验证本方法在材质质量、生成效率及一致性等方面的改进效果。  

（4）在系统实现方面，开发一个基于 Web 的交互平台，  
实现模型上传、材质生成、结果展示与简单交互等功能。  
用户可以通过浏览器上传白模或参考模型，  
在线完成材质生成，  
并实时预览结果，  
从而提升系统的易用性与实用价值。  

（5）在工程与应用层面，形成一套完整的三维材质生成与迁移流程，  
包括数据预处理、模型推理、结果输出等模块，  
为后续在游戏开发、数字内容创作及虚拟现实等领域的应用提供技术基础。  

综上，  
本课题预期实现从理论方法到实际系统的完整闭环，  
不仅验证深度学习在三维材质生成与迁移任务中的有效性，  
同时提升相关技术的可用性与应用价值。  
\subsection{研究计划}

本课题围绕三维物体材质生成与迁移任务，  
结合数据驱动方法与工程实现，  
计划按照如下阶段逐步开展研究工作：  

（1）前期调研与技术选型阶段。  
系统梳理三维材质生成与迁移领域的相关研究进展，  
重点关注基于扩散模型与多视角一致性优化的方法。  
同时，  
对Texverse数据集\cite{zhang2025texverse}及Trellis 2.0 \cite{xiang2025native}模型的结构与训练机制进行深入分析，  
明确模型改进方向与整体技术路线。  

（2）数据准备与预处理阶段。  
基于Texverse数据集\cite{zhang2025texverse}，  
构建适用于本任务的训练数据，  
包括带材质模型、对应白模以及必要的多视角渲染数据。  
对数据进行规范化处理，  
如统一分辨率、UV展开处理以及多视角图像对齐等，  
为后续模型训练提供高质量输入。  

（3）多视角生成模块构建。  
引入类似MV-Adapter\cite{huang2025mv}的多视角生成方法，  
将单视角输入扩展为多视角图像集合。  
通过预训练模型或微调策略，  
使生成的多视角图像在几何结构与语义上保持一致，  
为后续材质生成提供可靠的多视角条件信息。  

（4）模型改进与微调阶段。  
在Trellis 2.0\cite{xiang2025native}模型基础上进行改进，  
使其能够接收多视角输入作为条件信息。  
结合Texverse数据集对模型进行微调训练，  
重点优化多视角一致性、纹理细节表达以及几何结构适配能力。  
根据实验结果不断调整网络结构与损失函数设计。  

（5）实验验证与性能评估。  
设计定性与定量实验，  
对生成结果在纹理质量、多视角一致性及几何匹配程度等方面进行评估。  
通过与基线方法进行对比分析，  
验证所提出方法的有效性与优势。  

（6）系统实现与平台开发。  
基于Web技术开发交互平台，  
实现模型上传、材质生成、多视角预览等功能。  
将训练好的模型部署到后端服务中，  
支持用户在线调用，  
提高系统的实用性与可访问性。  

（7）总结与论文撰写阶段。  
对整个研究过程进行总结，  
整理实验结果与分析结论，  
完成毕业论文撰写，  
并对未来可能的改进方向与应用场景进行展望。  

通过上述研究计划，  
逐步实现从多视角生成到材质迁移的完整流程，  
提升三维材质生成的质量与稳定性。  

\subsection{进度安排}
\begin{table}[htbp]
\centering
\caption{三维物体材质生成与迁移课题进度安排}
\begin{tabularx}{\textwidth}{|c|c|X|}
\hline
任务 & 日期 & \multicolumn{1}{c|}{工作内容简述} \\
\hline
前期调研与技术选型 & 3.26-4.2 & 调研三维材质生成与迁移领域的相关研究，分析Texverse数据集与Trellis 2.0模型，确定技术路线与方法框架。 \\
\hline
数据准备与预处理 & 4.3-4.9 & 构建训练数据集，包括带材质模型、白模及多视角渲染数据，进行统一分辨率、UV处理和多视角对齐。 \\
\hline
多视角生成模块构建 & 4.10-4.16 & 引入多视角生成方法，实现单视角输入到多视角图像集合的生成，并确保几何与语义一致性。 \\
\hline
模型改进与微调 & 4.17-4.30 & 在Trellis 2.0基础上改进网络结构以支持多视角输入，结合Texverse数据集进行微调训练，优化多视角一致性与纹理细节。 \\
\hline
实验验证与性能评估 & 5.1-5.7 & 设计定性与定量实验，评估生成效果的纹理质量、多视角一致性及几何匹配，并与基线方法对比分析。 \\
\hline
系统实现与平台开发 & 5.8-5.13 & 开发基于Web的交互平台，实现模型上传、材质生成、多视角预览及简单交互功能，部署训练好的模型到后端服务。 \\
\hline
总结与论文撰写 & 5.14-5.16 & 整理实验结果与分析结论，完成论文初稿撰写，并总结整个研究过程与未来改进方向。 \\
\hline
\end{tabularx}
\end{table}
